\textbf{Решение задачи 2-4}

\vspace{1em}
\textbf{1. Исходные данные}

Из условия задачи известны следующие параметры:
\begin{itemize}
    \item Шумовая полоса БВЧ: $\Pi_{\text{ш}} = 1 \text{ МГц} = 10^6 \text{ Гц}$
    \item Температура согласованного генератора: $T = +20^\circ\text{C} = 293 \text{ К}$
    \item Постоянная составляющая напряжения на выходе АД без сигнала: $\bar{U}_{\text{ш}} = 0.2 \text{ В}$
    \item Мощность входного сигнала: $P_c = 5.07 \cdot 10^{-13} \text{ Вт}$
    \item Постоянная составляющая напряжения на выходе АД с сигналом: $\bar{U}_{\text{с+ш}} = 0.8 \text{ В}$
\end{itemize}
Постоянная Больцмана: $k = 1.38 \cdot 10^{-23} \text{ Вт/(Гц}\cdot\text{К)}$.

\vspace{1em}
\textbf{2. Определение отношения сигнал/шум на входе детектора}

Линейный амплитудный детектор (АД) выделяет огибающую входного процесса. При отсутствии полезного сигнала постоянная составляющая на выходе АД пропорциональна эффективному напряжению шума $\sigma$ на его входе:
\begin{equation*}
    \bar{U}_{\text{ш}} = K_{\text{д}} \sigma \sqrt{\frac{\pi}{2}}
\end{equation*}

При воздействии сильного гармонического сигнала ($U_m \gg \sigma$) постоянная составляющая на выходе АД приближенно равна амплитуде сигнала с учетом коэффициента передачи:
\begin{equation*}
    \bar{U}_{\text{с+ш}} \approx K_{\text{д}} U_m
\end{equation*}

Найдем отношение амплитуды сигнала к СКО шума ($a = U_m/\sigma$):
\begin{equation*}
    \frac{\bar{U}_{\text{с+ш}}}{\bar{U}_{\text{ш}}} = \frac{K_{\text{д}} U_m}{K_{\text{д}} \sigma \sqrt{\frac{\pi}{2}}} = \frac{U_m}{\sigma} \sqrt{\frac{2}{\pi}} = a \sqrt{\frac{2}{\pi}}
\end{equation*}
\begin{equation*}
    a = \frac{\bar{U}_{\text{с+ш}}}{\bar{U}_{\text{ш}}} \sqrt{\frac{\pi}{2}} = \frac{0.8}{0.2} \sqrt{\frac{\pi}{2}} = 4 \sqrt{\frac{\pi}{2}} \approx 4 \cdot 1.2533 = 5.013
\end{equation*}
Поскольку $a \approx 5 > 3$, использование асимптотики для сильного сигнала полностью оправдано.

Отношение сигнал/шум по мощности на входе АД (и на входе приемника) равно:
\begin{equation*}
    q^2 = \frac{P_c}{P_{\text{ш}}} = \frac{U_m^2 / 2}{\sigma^2} = \frac{1}{2} a^2
\end{equation*}
Подставим найденное значение $a$:
\begin{equation*}
    q^2 = \frac{1}{2} \left( 4 \sqrt{\frac{\pi}{2}} \right)^2 = \frac{1}{2} \cdot 16 \cdot \frac{\pi}{2} = 4\pi \approx 12.566
\end{equation*}

\vspace{1em}
\textbf{3. Расчет коэффициента шума БВЧ}

Выразим мощность собственных шумов приемника, приведенную ко входу:
\begin{equation*}
    P_{\text{ш}} = \frac{P_c}{4\pi} = \frac{5.07 \cdot 10^{-13}}{12.566} \approx 4.035 \cdot 10^{-14} \text{ Вт}
\end{equation*}

С другой стороны, мощность шума связана с шумовой температурой источника и коэффициентом шума БВЧ формулой Найквиста:
\begin{equation*}
    P_{\text{ш}} = F_{\text{БВЧ}} \cdot k T \Pi_{\text{ш}}
\end{equation*}

Рассчитаем мощность тепловых шумов согласованного генератора при заданной температуре $T = 293 \text{ К}$:
\begin{equation*}
    P_{\text{т}} = k T \Pi_{\text{ш}} = 1.38 \cdot 10^{-23} \cdot 293 \cdot 10^6 \approx 4.043 \cdot 10^{-15} \text{ Вт}
\end{equation*}

Найдем искомый коэффициент шума БВЧ (в разах):
\begin{equation*}
    F_{\text{БВЧ}} = \frac{P_{\text{ш}}}{P_{\text{т}}} = \frac{4.035 \cdot 10^{-14}}{4.043 \cdot 10^{-15}} \approx 9.98 \approx 10 \text{ раз}
\end{equation*}

Переведем полученное значение в децибелы:
\begin{equation*}
    F_{\text{БВЧ (дБ)}} = 10 \log_{10}(10) = 10 \text{ дБ}
\end{equation*}

\textbf{Ответ:} Коэффициент шума БВЧ равен $10$ (или $10 \text{ дБ}$).